\section{Complements}
The notion of a complement requires the idea of an \emph{universal set}. The way I think about it is thus. Let $P$ be the set of prime numbers, then if asked to give a number outside of $P$, one can point out to examples of composite numbers such as $4,6, \cdots$. So there is an underlying assumption that $P$ is part of a bigger set which is the set of natural numbers, i.e. $P\subseteq\N$. In this case we call $\N$ the \emph{universe} or the \emph{universal set.}
\begin{defn}
    Let $A$ be a set with a universal set $U$. The complement of $A$, denoted $A^{\complement}$ is the set $A^{\complement}=U-A$.
\end{defn}

So, if $P$ is the set of prime numbers, then: 
$$
P^{\complement}=\N-P=\{ 1,4,6,8,9,10\dots \}.
$$\footnote{If $0\in\N$, then the set contains $0$ as well.}
and so we can say that $P^{\complement}$ is the set of all \emph{composite numbers and $1$}. 

\subsection{Exercises for Section 6}
\begin{exercise}\label{ex:6.1}
    Let $A=\{ 4,3,6,7,1,9 \}, B=\{ 5,6,8,4 \}$ have universal set $U=\{ 0,1,2,\dots,10 \}$. Find: 
    \begin{enumerate}
        \item $A^{\complement}$
        \item $B^{\complement}$
        \item $A\cap A^{\complement}$
        \item $A\cup A^{\complement}$
        \item $A-A^{\complement}$
        \item $A-B^{\complement}$
        \item $A^{\complement}-B^{\complement}$
        \item $A^{\complement}\cap B$
        \item $(A^{\complement}-B)^{\complement}$
    \end{enumerate}
\end{exercise}\hyperref[sol:6.1]{[Go to Solution]}

\begin{exercise}\label{ex:6.2}
    Let $A=\{ 0,2,4,6,8 \}, B=\{ 1,3,5,7 \}$ have universal set $U=\{ 0,1,2,\dots,8 \}$. Find: 
    \begin{enumerate}
        \item $A^{\complement}$
        \item $B^{\complement}$
        \item $A\cap A^{\complement}$
        \item $A\cup A^{\complement}$
        \item $A-A^{\complement}$
        \item $(A\cup B)^{\complement}$
        \item $A^{\complement}\cap B^{\complement}$
        \item $(A\cap B)^{\complement}$
        \item $A^{\complement}\times B$
    \end{enumerate}
\end{exercise}\hyperref[sol:6.2]{[Go to Solution]}

\begin{exercise}\label{ex:6.3}
    Sketch the set $X=[1,3]\times[1,2]$ on the plane $\R^{2}$. On separate drawings, shade in the sets $X^{\complement}$ and $X^{\complement}\cap([0,2]\times[0,3])$.
\end{exercise}\hyperref[sol:6.3]{[Go to Solution]}

\begin{exercise}\label{ex:6.4}
    Sketch the set $X=[-1,3]\times[0,2]$ on the plane $\R^{2}$. On separate drawings, shade in the sets $X^{\complement}$ and $X^{\complement}\cap([-2,4]\times[-1,3])$.
\end{exercise}\hyperref[sol:6.4]{[Go to Solution]}

\begin{exercise}\label{ex:6.5}
    Sketch the set $X=\{ (x,y)\in\R^{2}:1\le x^{2}+y^{2}\le 4 \}$ on the plane $\R^{2}$. On separate drawings, shade in the sets $X^{\complement}$
\end{exercise}\hyperref[sol:6.5]{[Go to Solution]}

\begin{exercise}\label{ex:6.6}
    Sketch the set $X=\{ (x,y)\in\R^{2}:y<x^{2}\}$ on the plane $\R^{2}$. Shade in the sets $X^{\complement}$
\end{exercise}\hyperref[sol:6.6]{[Go to Solution]}